\chapter{实射影空间}

\section{实射影空间的拓扑}

记号约定:$\Delta_{n}$表示$\mathbb{P}_{n}\left(\R\right)$上的等价关系$\forall\mathbf{x},\mathbf{y}\in\R^{n}\left(x\Delta_{n}\mathbf{y}\iff \mathbf{y}\in\R^{\times}\mathbf{x}\right)$.

\begin{definition}{实射影空间,实射影直线,实射影平面}{}
	给$\mathbb{P}_{n}\left(\R\right)$装备商拓扑,所得拓扑空间记作$\RP^{n}$,称为$n$维实射影空间,$\RP^{1},\RP^{2}$分别称为实射影直线和实射影平面.
\end{definition}

\begin{proposition}{$\RP^{n}$的Hausdorff性}{}
	$\RP^{n}$是Hausdorff空间.
\end{proposition}

\begin{proposition}{$\RP^{n}$的紧性和连通性}{}
	$\RP^{n}$紧且连通.
\end{proposition}

\begin{proposition}{$\RP^{n}$的球面模型}{}
	如果$\Delta_{n}^{\prime}$是$\Delta_{n}$在$\S_{n}$上诱导的等价关系,那么$\RP^{n}$同胚于$\S_{n}/\Delta_{n}^{\prime}$.
\end{proposition}

\begin{proposition}{$\RP^{n}$的球体模型}{}
	如果$\Delta_{n}^{\prime}$是$\Delta_{n}$在$\B_{n}$上诱导的等价关系,那么$\RP^{n}$同胚于$\B_{n}/\Delta_{n}^{\prime}$.
\end{proposition}